\section{Implementación en CODESYS}

\begin{frame}{De Stateflow a un PLC (IEC 61131)}
  Para llevar el control a hardware real se traduce la lógica a CODESYS. La consigna pide
  \textbf{dos métodos} y compararlos:
  \begin{columns}[T]
    \begin{column}{0.5\textwidth}
      \begin{block}{1 --- Codificación manual}
        Estructura en \textbf{SFC} (Sequential Function Charts) + cálculos en \textbf{ST}
        (Structured Text).
      \end{block}
    \end{column}
    \begin{column}{0.5\textwidth}
      \begin{block}{2 --- Generación automática}
        \textbf{PLC Coder} de Simulink: todo a \textbf{ST} desde los Statecharts.
      \end{block}
    \end{column}
  \end{columns}
  \centering
  \footnotesize Niveles 0 y 1 como POUs cíclicos en paralelo; Nivel 2 en Simulink (simil ST).
\end{frame}

\begin{frame}{Codificación manual: SFC + ST}
  \begin{columns}[T]
    \begin{column}{0.32\textwidth}
      \centering
      \includegraphics[width=\linewidth]{codesys/arbol_codesys.png}\\
      \footnotesize Árbol del proyecto
    \end{column}
    \begin{column}{0.68\textwidth}
      \centering
      \includegraphics[width=\linewidth]{codesys/codesys_supervisor11.png}\\
      \footnotesize Supervisor en SFC (mismos estados y transiciones que Stateflow)
    \end{column}
  \end{columns}
\end{frame}

\begin{frame}{Manual (SFC) vs. autogenerado (ST)}
  \begin{columns}[T]
    \begin{column}{0.5\textwidth}
      \begin{itemize}
        \item Rendimiento y carga de CPU \textbf{comparables}.
        \item El ST autogenerado es \textbf{monolítico} (un \texttt{CASE} global + \texttt{IF}
              anidados) $\Rightarrow$ poca trazabilidad para diagnóstico en planta.
        \item Se eligió \textbf{SFC manual} para la arquitectura (visual) y ST sólo para los
              cálculos.
      \end{itemize}
    \end{column}
    \begin{column}{0.5\textwidth}
      \centering
      \includegraphics[width=\linewidth]{codesys/st_generado.png}\\
      \footnotesize Fragmento de ST autogenerado (PLC Coder)
    \end{column}
  \end{columns}
\end{frame}

\begin{frame}{Co-simulación CODESYS $\leftrightarrow$ Simulink (OPC-UA)}
  \begin{columns}[T]
    \begin{column}{0.45\textwidth}
      \begin{itemize}
        \item PLC virtual (CODESYS) + planta (Simulink) comunicados por \textbf{OPC-UA}.
        \item Software-in-the-Loop: el controlador real maneja la planta simulada.
        \item Se verifica que responde igual que el modelo original.
      \end{itemize}
    \end{column}
    \begin{column}{0.55\textwidth}
      \centering
      \includegraphics[width=\linewidth]{codesys/codesys_interfaz.png}
    \end{column}
  \end{columns}
\end{frame}
